\BourbakiExerciseGroup

本节中的所有域，除非另有明确说明，均为交换域。

\begin{enumerate}
\def\labelenumi{\arabic{enumi})}
\item
  设 u 是有限维向量空间 E 的一个自同态，并设 V 是 E 中在 u 下不变的子空间。证明 u 在 V 上的限制的特征多项式整除 u 的特征多项式。由此得出 Cayley-Hamilton 定理的一个初等证明（化归到 \(E_{u}\) 是循环的情形，并利用 VII 第 30 页的公式 (2)）。
\item
  \begin{enumerate}
  \def\labelenumii{\alph{enumii})}
  \tightlist
  \item
    证明域 K 上的每个方阵都相似于其转置（化归到 Jordan 矩阵的情形）。
  \end{enumerate}
\end{enumerate}

\begin{enumerate}
\def\labelenumi{\alph{enumi})}
\setcounter{enumi}{1}
\item
  证明在有理整数环 \(\mathbf{Z}\) 上，矩阵 \(\begin{pmatrix} 8 & 2 \\ 0 & 1 \end{pmatrix}\) 不相似于其转置。
\item
  设 \(u\) 是域 \(\mathbf{K}\) 上有限维向量空间的一个自同态。则每个特征值 \(\lambda\) （属于 \(u\) ）也是转置自同态 \(^t u\) （属于 \(E^*\) ）的特征值。此外，若 \(\lambda\) 与 \(\mu\) 是 \(u\) 的两个不同的特征值，\(x\) 是 \(u\) 的对应于 \(\lambda\) 的特征向量，且 \(y^*\) 是 \(^t u\) 的对应于 \(\mu\) 的特征向量，则 \(\langle x, y^* \rangle = 0\) 。
\end{enumerate}

\begin{enumerate}
\def\labelenumi{\arabic{enumi})}
\setcounter{enumi}{2}
\item
  设 u 是有限维向量空间 E 的一个自同态，令 \(\lambda\) 为 u 的一个特征值，并令 v 为自同态 \(u - \lambda\) 。证明 E 是 \(v^{-1}(0)\) 与 \(v(E)\) 的直和，当且仅当 \(\lambda\) 是 u 的极小多项式的单根。
\item
  通过向基域添加一个超越元，证明 VII 第 36 页命题 10 中的公式 (6) 可以从公式 (8) 推出。通过将多项式 q 分解为线性因式来直接证明后者。
\item
  \begin{enumerate}
  \def\labelenumii{\alph{enumii})}
  \tightlist
  \item
    设 K 是一个代数闭域，A 是 K 上一个 n 阶方阵，呈对角分块 \(\text{diag}(A_i)\) （由诸下三角矩阵 \(A_i\) 组成）的形式，使得 \(A_i\) 的对角元都等于 K 的同一个元素 \(\alpha_i\) ( \(1 \leqslant i \leqslant r\) )，且 \(\alpha_i \neq \alpha_j\) 对 \(i \neq j\) 成立。证明每个与 A 可交换的方阵 B 都是对角分块 \(\text{diag}(B_i)\) ，其中 \(B_i\) 与 \(A_i\) 同阶，对 \(1 \leqslant i \leqslant r\) 。
  \end{enumerate}
\end{enumerate}

\begin{enumerate}
\def\labelenumi{\alph{enumi})}
\setcounter{enumi}{1}
\tightlist
\item
  设 M 是 K 上一些 n 阶可交换方阵组成的一个集合。证明通过对 M 中所有矩阵施加相同的相似变换，可使它们都成为下三角矩阵的分块对角矩阵 \(\text{diag}(X_i)\) ，其中 \(X_i\) 的个数 r 以及 \(m_i\) （即 \(X_i\) 对每个 i 的阶）对 M 中每个矩阵都相同，且 \(X_i\) 的所有特征值都相等。（利用 a) 化归到 r = 1 的情形，然后利用 VII 第 41 页引理 3 证明在此情形 M 中所有矩阵都有非零的公共特征向量，再对 n 作归纳完成证明。）
\end{enumerate}

\begin{enumerate}
\def\labelenumi{\arabic{enumi})}
\setcounter{enumi}{5}
\item
  交换域 K 上两个由 n 阶方阵组成的对 \((A_{1}, A_{2})\) 与 \((B_{1}, B_{2})\) 称为等价的，如果存在 K 上两个可逆方阵 P 和 Q，使得 \(A_{1} = PB_{1}Q\) 且 \(A_{2} = PB_{2}Q\) 。若 \(A_{1}\) 与 \(B_{1}\) 可逆，证明这两个对等价当且仅当矩阵 \(XA_{1} + A_{2}\) 与 \(XB_{1} + B_{2}\) 在环 K{[}X{]} 上等价。
\item
  设 \(A\) 是一个 \(m\) 行 \(n\) 列、秩为 \(r\) 的矩阵，定义在域 \(K\) 上。若 \(P\) 是一个 \(m\) 阶方阵而 \(Q\) 是一个 \(n\) 阶方阵（均在 \(K\) 上），使得 \(PAQ = A\) ，证明存在一个矩阵 \(P'\) 与 \(P\) 相似，以及一个矩阵 \(Q'\) 与 \(Q\) 相似，使得
\end{enumerate}

\[
P ^ {\prime} = \left( \begin{array}{c c} R & S \\ 0 & U \end{array} \right), \quad Q ^ {\prime} = \left( \begin{array}{c c} R ^ {- 1} & 0 \\ S ^ {\prime} & U ^ {\prime} \end{array} \right)
\]

其中R是一个r阶可逆方阵，U（相应地 \(U'\) ) 是 m - r 阶（相应地，n - r 阶）方阵，而 S（相应地， \(S'\) ）是一个具有 r 行和 m - r 列（相应地，n - r 行和 m 列）的矩阵。

\begin{enumerate}
\def\labelenumi{\arabic{enumi})}
\setcounter{enumi}{7}
\item
  设 E 是域 K 上有限维为 n 的向量空间，u 是 E 的一个自同态，V 是 E 中在 u 下不变的子空间，\(u|V\) 是 u 在 V 上的限制。若 \(f_{i}\) ( \(1 \leqslant i \leqslant n\) )（相应地 \(g_{j}\) ( \(1 \leqslant j \leqslant m\) )）是 u（相应地 \(u|V\) ）的相似不变量，证明 \(g_{j}\) 整除 \(f_{j+n-m}\) ，对 \(1 \leqslant j \leqslant m\) （参见 VII，第 64 页，习题 8，d)）。陈述并证明 u 诱导的 E/V 的自同态的类似性质。
\item
  设 \(A\) 是一个 \(n\) 阶方阵，\(B\) 是一个 \(m\) 阶方阵，\(C\) 是一个 \(m\) 行 \(n\) 列的矩阵，定义在域 \(K\) 上。把方程 \(X A = B X + C\) （其中 \(X\) 是一个 \(m\) 行 \(n\) 列的未知矩阵，定义在 \(K\) 上）化为多项式环 \(K[Y]\) 中的线性同余式方程组（对所有 \(f \in K[Y]\) ，注意 \(X f(A) = f(B)X + H(f)\) 对某个完全确定的矩阵 \(H(f)\) （ \(m\) 行 \(n\) 列）成立；令 \(\oplus_{i} E_i\) 为模 \(E\) （与 \(A\) 相伴，见第 1 节）分解为
\end{enumerate}

循环模的直和分解，这些循环模的零化理想即是 A 的异于 1 的相似不变量 \(f_{i}\) （命题 2），并令 \(a_{i}\) 为诸 \(E_{i}\) 的生成元；同样地定义子模 \(F_{j}\) （由 \(b_{j}\) 生成）以及对应于 B 的多项式 \(g_{j}\) ；写出 \(Xf_{i}(A)a_{i}=0\) ；然后把 \(Xa_{i}\) 与 \(H(f_{i})a_{i}\) 按诸 \(F_{j}\) 分解，并证明这样从方程 \(Xf_{i}(A)a_{i}=0\) 得到的条件对解的存在性是必要且充分的）。

\begin{enumerate}
\def\labelenumi{\arabic{enumi})}
\setcounter{enumi}{9}
\tightlist
\item
  \begin{enumerate}
  \def\labelenumii{\alph{enumii})}
  \tightlist
  \item
    证明：在任意域 K 上，矩阵
  \end{enumerate}
\end{enumerate}

\[
\left( \begin{array}{c c c c c c} \lambda & \alpha & \beta_ {13} & \beta_ {14} & ... & \beta_ {1n} \\ 0 & \lambda & \alpha & \beta_ {24} & ... & \beta_ {2n} \\ 0 & 0 & \lambda & \alpha & ... & \beta_ {3n} \\ \cdots & \cdots & \cdots & \cdots & \cdots & \cdots \\ 0 & 0 & 0 & 0 & ... & \alpha \\ 0 & 0 & 0 & 0 & ... & \lambda \end{array} \right)
\]

相似于若尔当矩阵 \(U_{n,\lambda}\) ，只要 \(\alpha \neq 0\) 。

\begin{enumerate}
\def\labelenumi{\alph{enumi})}
\setcounter{enumi}{1}
\item
  若 \(\lambda \neq 0\) 且 \(\mathbf{K}\) 的特征为 0，证明 \((U_{n,\lambda})^m\) 相似于若尔当矩阵 \(U_{n,\lambda^m}\) ，对每个整数 \(m\) （正的或负的）。
\item
  证明 \((U_{n,0})^m = 0\) 对 \(m \geqslant n\) 。对 \(0 < m < n\) ，设 \(n - 1 = km + q\) ，其中 \(0 \leqslant q < m\) ；证明 \((U_{n,0})^m\) 有 \(q + 1\) 个等于 \(\mathbf{X}^{k+1}\) 的相似不变量，以及 \(m - q - 1\) 个等于 \(\mathbf{X}^k\) 的相似不变量（把 \(\mathbf{K}^n\) 的典范基按适当的顺序排列）。
\item
  类似地，确定 \(f(U_{n,\lambda})\) （ \(f \in \mathbf{K}[\mathbf{X}]\) ）的相似不变量。
\item
  对特征为 \(p \neq 0\) 的域 K 讨论同样的问题。
\end{enumerate}

\begin{enumerate}
\def\labelenumi{\arabic{enumi})}
\setcounter{enumi}{10}
\item
  设 K 为一个代数闭域，设 A 为 K 上的一个可逆矩阵，且 m \textgreater{} 0 是一个整数。证明：如果 m 不是 K 的特征的倍数，则存在一个多项式 \(g(\mathbf{X}) \in \mathbf{K}[\mathbf{X}]\) 使得 \((g(A))^{m} = A\) （利用凯莱-哈密顿定理及 VII，第 54 页，习题 28）。把方程推广到 \(f(U) = A\) ( \(f \in K[X]\) ，U 为 K 上未知的 \(n \times n\) 矩阵）。
\item
  \begin{enumerate}
  \def\labelenumii{\alph{enumii})}
  \tightlist
  \item
    设 \(A\) 与 \(B\) 是两个 \(n\) 阶方阵，定义在交换域 \(\mathbf{K}\) 上，并设 \(f_i\) 与 \(g_i\) ( \(1 \leqslant i \leqslant n\) ) 分别是它们的相似不变量。证明在 \(\mathbf{K}\) 上由方阵 \(U\) （ \(n\) 阶且满足 \(UA = BU\) ）组成的向量空间同构于如下定义的空间 \(\mathbf{M}_0 = \mathbf{M} / \mathbf{N}\) ：\(\mathbf{M}\) 是方阵 \((u_{ij}(\mathbf{X}))\) 组成的空间，这些方阵是 \(n\) 阶的、定义在 \(\mathbf{K}[\mathbf{X}]\) 上，且满足 \(u_{ij}(\mathbf{X}) f_j(\mathbf{X})\) 是 \(g_i(\mathbf{X})\) 的倍式，对每对指标 \(i, j\) ；而 \(\mathbf{N}\) 是 \(\mathbf{M}\) 的子空间，由满足 \(u_{ij}(\mathbf{X})\) 是 \(g_i(\mathbf{X})\) 的倍式（对每对指标 \(i, j\) ）的那些矩阵组成。（把 \(U, A\) 与 \(B\) 分别看作自同态 \(u, v\) 与 \(w\) （属于 \(E = K^n\) ）的矩阵，并注意 \(u\) 是一个 \(K[X]\) -线性映射，从 \(E_v\) 到 \(E_w\) ；然后把 \(E_v\) 与 \(E_w\) 表示为商模。）
  \end{enumerate}
\end{enumerate}

\begin{enumerate}
\def\labelenumi{\alph{enumi})}
\setcounter{enumi}{1}
\item
  当 B = A 时，满足 UA = AU 的矩阵 U 构成一个环 P；此时 M 是一个方阵环，N 是 M 的一个双边理想，而 P 同构于商环 M/N。证明 P 与环 K{[}A{]} （由 \(I_{n}\) 和 A 生成）相同，当且仅当 A 的极小多项式等于它的特征多项式。
\item
  对任意的 \(A\) 与 \(B\) ，证明 \(\mathbf{M}_0\) 的维数等于 \(\sum_{i,j} m_{ij}\) ，其中
\end{enumerate}

\(m_{ij}\) 是 \(f_i\) 与 \(g_j\) 的最大公因子的次数（与 \(a\) 的方法相同）。

\begin{enumerate}
\def\labelenumi{\alph{enumi})}
\setcounter{enumi}{3}
\tightlist
\item
  证明矩阵 \(U \in \mathbf{M}_0\) 的秩的最大值等于 \(\sum_{k} d_k\) ，其中
\end{enumerate}

\(d_{k}\) 是 \(f_{k}\) 与 \(\pmb{g}_{k}\) 的最大公因子的次数（利用习题 8）。

\begin{enumerate}
\def\labelenumi{\alph{enumi})}
\setcounter{enumi}{4}
\tightlist
\item
  把 a)、c) 和 d) 的结果推广到如下情形：A 是 n 阶方阵，B 是 m 阶方阵，而 U 是 m 行 n 列的矩阵。
\end{enumerate}

\begin{enumerate}
\def\labelenumi{\arabic{enumi})}
\setcounter{enumi}{12}
\tightlist
\item
  设 \((\mathbf{Z}_{ij})\) ( \(1 \leqslant i, j \leqslant n\) ) 是 \(n^2\) 个不定元，定义在域 \(\mathbf{K}_0\) 上。设 \(\mathbf{A}\) 为多项式环 \(\mathbf{K}_0[\mathbf{Z}_{11}, \ldots, \mathbf{Z}_{nn}]\) ，并设 \(\mathbf{K} = \mathbf{K}_0(\mathbf{Z}_{ij})\) 为其分式域；令 \(\chi_U\) 为矩阵 \(U = (\mathbf{Z}_{ij})\) 在 \(\mathbf{K}\) 上的特征多项式。
\end{enumerate}

\begin{enumerate}
\def\labelenumi{\alph{enumi})}
\item
  证明 \(\chi_{U}\) 是一个可分的不可约多项式（证明若 \(\chi_{U}\) 在 K 上可约，则它将是 \(K_{0}[X, Z_{ij}]\) 中两个关于 X 的次数 \textgreater0 的多项式之积；由此推出 K 上每个 n 阶方阵的特征多项式都会是可约的；然后从系数在 \(K_{0}\) 上代数无关的多项式的例子导出矛盾；用类似的方法证明可分性）。
\item
  证明 U 相似于一个每个元素都等于 0、1 或多项式 \(\chi_{U}\) 的某个系数的矩阵（参见 VII，第 30 页，公式 (2)）。由此推出：若 \(K_{0}\) 是一个无限域，p 是 A 中一个多项式，使得 \(p(u_{ij}^{\prime}) = p(u_{ij}^{\prime\prime})\) 对每对相似的矩阵 \((u_{ij}^{\prime})\) 与 \((u_{ij}^{\prime\prime})\) （在 \(K_{0}\) 上）成立，则 p 属于由 \(\chi_{U}\) 的系数生成的 A 的子环。对有理函数 \(r \in K = K_{0}(Z_{ij})\) 陈述并证明类似的结果。
\end{enumerate}

\begin{enumerate}
\def\labelenumi{\arabic{enumi})}
\setcounter{enumi}{13}
\tightlist
\item
  设 K 是特征为 p 的非完全域，设 a 是 \(K - K^{p}\) 的一个元素，并设 L 是域 \(\mathrm{K}(a^{1/p})\) 。考虑 K-向量空间 \(E = L \otimes_{K} L\) ；设 u（相应地 \(u'\) ）是 E 的 K-自同态，由在代数 E 中乘以 \(a^{1/p} \otimes 1\) （相应地 \(1 \otimes a^{1/p}\) ）给出。证明 u 与 \(u'\) 都是半单的，u 与 \(u'\) 可交换，且 \(u - u'\) 是非零幂零的。
\end{enumerate}

